%% EXHALE: what the energy-stable WENO (ESWENO) schemes offer, and what of
%% them the code already uses. Written 2026-09-12 from the three published
%% papers in ../references/ (Yamaleev & Carpenter 2009a,b; Fisher et al.
%% 2011) and from src/modules/states/Reconstruction.f90.
%% Author: Kwang-Il Seon
\documentclass[english,a4paper]{my_memo}
\synctex=1
\usepackage{booktabs}
\usepackage{amsmath}
\usepackage[authoryear,round]{natbib}
\usepackage{babel}
\emergencystretch=3em

\title{Energy-stable WENO: what the three papers establish and what EXHALE uses}
\author{Kwang-Il Seon}
\date{Last updated: \docmoddate}

\begin{document}
\maketitle

\begin{abstract}
The energy-stable WENO (ESWENO) schemes of \citet{Yamaleev2009a,Yamaleev2009b}
and their finite-domain closures \citep{Fisher2011} add to a conventional WENO
flux a nonlinear artificial dissipation term built from the WENO weights
themselves, and prove that the resulting finite-difference operator is
$L_2$-stable for linear hyperbolic equations with continuous or discontinuous
data, with no total-variation assumption and no time-step restriction beyond
the usual one. The same papers introduce new weight functions that hold the
design order at smooth extrema with any number of vanishing derivatives, and
the 2011 paper supplies summation-by-parts boundary closures so the estimate
holds on a finite domain. EXHALE's WENO3 reconstruction already uses the
2009a weight functions and their $\epsilon=O(\Delta\xi^2)$ floor; it does not
compute the dissipation term, so it inherits the accuracy result and not the
energy estimate. Three properties of the code keep the estimate from carrying
over as published: a finite-volume reconstruction of the primitive state on a
stretched grid, a component-wise rather than characteristic reconstruction,
and the nonlinear Euler system, for which the papers themselves give no
proof. This memo records what each paper proves, at which equation, and what
would have to change for the estimate to apply.
\end{abstract}

\section{The question}

The question asked on 2026-09-12 was what ESWENO gains over conventional
WENO. The three papers are in \texttt{../references/} (the files
\texttt{Yamaleev\_2009JCP\_228\_3025.pdf}, \texttt{Yamaleev\_2009JCP\_228\_4248.pdf}
and \texttt{Fisher\_2011JCP\_230\_3727.pdf}) and every statement below was read
in those published versions; equation numbers are theirs. What the code does
is read from \texttt{src/modules/states/Reconstruction.f90}, whose WENO3
branch carries its own account of the same papers.

\section{The third-order scheme \citep{Yamaleev2009a}}

\subsection{The construction}

The scalar wave equation $u_t + f_x = 0$, $f = au$, $a \ge 0$, is
discretized on a uniform periodic grid as (their Eq.~9)
\begin{equation}
  \frac{\partial \mathbf{u}}{\partial t} + P^{-1} Q \mathbf{f}
  = -\,P^{-1} D_1^{T} S D_1 \mathbf{f},
  \label{eq:esweno}
\end{equation}
where $P^{-1}Q$ is the skew-symmetric part of the conventional third-order
WENO derivative matrix, $D_1$ the two-point difference operator, and $S$ a
symmetric matrix that is a nonlinear function of the solution through the
WENO weights. The left side is the WENO scheme; the right side is the
addition. Multiplying by $\mathbf{u}^{T}P$ gives (their Eq.~11)
\begin{equation}
  \frac{d}{dt}\|\mathbf{u}\|_P^2 + a\,\mathbf{u}^{T}(Q+Q^{T})\mathbf{u}
  = -\,a\,(D_1\mathbf{u})^{T}(S+S^{T})D_1\mathbf{u},
\end{equation}
so the energy cannot grow if $Q+Q^{T}$ vanishes in the periodic interior and
$S$ is positive semidefinite. Their $S$ (Eqs.~35 to 37) is
\begin{equation}
  S = S_1 + D_1^{T} S_2 D_1,\qquad
  s^{1}_{jj} = \mu_j - \tfrac{1}{8}\big(w^{1}_{j+1/2} - w^{1}_{j-1/2}\big),\qquad
  s^{2}_{jj} = \tfrac{1}{4} w^{1}_{j-1/2},
\end{equation}
\begin{equation}
  \mu_j = \tfrac{1}{8}\sqrt{\big(w^{1}_{j+1/2} - w^{1}_{j-1/2}\big)^2 + \delta^2},
\end{equation}
with $w^{1}$ the weight of the upwind stencil and $\delta$ a small positive
parameter that keeps $\mu_j$ a $C^1$ function of the solution. The proof of
positive semidefiniteness (their Section~4.1) uses only two properties of
the weights, their Eq.~23: $0 \le w^{r} \le 1$ and $w^{0} + w^{1} = 1$. It
holds for any weight functions with those properties, continuous or
discontinuous data alike (their Remark~6), and fails if a weight is negative
(Remark~7). The stencil is the five points of conventional WENO3 (Remark~4).

The dissipation term is third order on smooth data, so the scheme keeps the
design order there; it is what makes the operator's symmetric part
negative semidefinite. The papers state as the distinctive feature that the
symmetric part of the conventional WENO operator can have positive
eigenvalues, so conventional WENO ``may become locally unstable''
(their Section~7), and that no energy estimate is available for it.

\subsection{The new weight functions}

Independently of the dissipation term, the paper replaces the
\citet{Jiang1996} weights $\alpha_r = d_r/(\epsilon + \beta_r)^2$ by
(their Eqs.~18, 21, 22)
\begin{equation}
  w^{r} = \frac{\alpha_r}{\alpha_0 + \alpha_1},\qquad
  \alpha_r = d_r\Big(1 + \frac{\tau}{\epsilon + \beta_r}\Big),\qquad
  \tau = (u_{j+1} - 2u_j + u_{j-1})^2,
  \label{eq:weights}
\end{equation}
with $\beta_0 = (u_{j+1}-u_j)^2$, $\beta_1 = (u_j - u_{j-1})^2$,
$d_0 = 2/3$, $d_1 = 1/3$. On smooth data $w^{r} = d_r + O(\Delta x^2)$ (their
Eq.~24) with no assumption on the number of vanishing derivatives, provided
$\epsilon \ge O(\Delta x^2)$ (Eq.~51); with $\epsilon = 0$ the ratio
$\tau/\beta_r$ is $O(1)$ at a smooth extremum and both WENO and ESWENO drop
to second order there (Eq.~57), which is what the Jiang--Shu weights do on
practical grids (Eq.~58). Nullifying the stencil astride a discontinuity
requires $\epsilon \le O(\Delta x^2)$ (Eq.~61), so $\epsilon = O(\Delta x^2)$
(Eq.~62), made scale-invariant with the computational spacing
$\Delta\xi = 1/J$ (Eq.~64) and given the solution scale by
\begin{equation}
  \epsilon = \max\big(\|u_0^2\|,\ \|(u_0)_\xi^2\|\big)\,\Delta\xi^2,\qquad
  \delta = \Delta\xi^2
\end{equation}
(Eqs.~65, 66; $L_1$ norms of the initial condition, evaluated once). The
paper says explicitly that ``these weights can also be used with the
conventional third-order WENO scheme'' (Section~3.1): the accuracy result
and the energy estimate are separable, and only the latter needs the
dissipation term.

\subsection{Scope of the proof}

The estimate is proven for the scalar equation and, in Section~5, for a
system of constant-coefficient linear hyperbolic equations in
\emph{characteristic} form, ``while it is not clear how to obtain a similar
energy estimate when the system of equations is discretized in the
component-wise fashion.'' For the Euler equations the paper is explicit
(Section~6.2): ``No stability proof is currently available for the ESWENO
scheme for the nonlinear Euler equations and the above characteristic
decomposition which is done locally at each point.'' The nonlinear tests
(quasi-one-dimensional nozzle, subsonic and transonic; time-dependent
problems with strong shocks and contacts) are reported as stable and
essentially non-oscillatory, with a local Roe-averaged characteristic
decomposition and global Lax--Friedrichs splitting, but as observations.

\section{Arbitrary order \citep{Yamaleev2009b}}

The second paper turns the construction into a procedure for any order:
decompose the symmetric part of the WENO operator into a sum
$\sum_{i=0}^{s}[D_1^{i}]\Lambda_i[D_1^{i}]^{T}$ of nested two-point
differences with diagonal $\Lambda_i$ (their Lemma~1 and Appendix),
replace each diagonal entry by its smoothly positive part
$\tfrac12\big(\sqrt{\lambda^2+\delta_i^2} - \lambda\big)$ (Eq.~23), and add
the resulting positive semidefinite operator (Eq.~24) as the dissipation
term. Fifth-order upwind-biased schemes as a one-parameter family, a
sixth-order central scheme, and schemes to eighth order are given, with
weight functions of the form \eqref{eq:weights} and constraints on their
parameters that keep design order with any number of vanishing derivatives
(their Eq.~58 and the floor of Eq.~79 are the higher-order counterparts of
the 2009a forms).

The measured cost (their Table~1, linear wave equation, 101 points, five
periods): the third-, fifth- and sixth-order ESWENO schemes take 15 to 20
percent more CPU time than the WENO schemes they are built on, the
fourth-order one somewhat more; no attempt at optimization was made. The
stencil width is unchanged.

\section{Boundary closures \citep{Fisher2011}}

The 2009 estimates are for periodic domains. The 2011 paper builds
closures for the fourth-order ESWENO scheme that satisfy the
summation-by-parts convention, $D = P^{-1}[Q + R]$ with
$Q + Q^{T} = \mathrm{diag}[-1, 0, \ldots, 0, 1]$ and $R$ symmetric positive
semidefinite (their Eqs.~4, 5), so that with a simultaneous-approximation-term
boundary condition of strength $\sigma \le -1/2$ the semi-discrete energy is
bounded (Theorem~1). The finding that makes it possible is that the flux
points near a boundary must be \emph{nonuniform} and derived from the norm
$P$ (their Section~3.3), which is what lets accuracy, the SBP convention and
the WENO stencil biasing hold at once. Second-order closures stable in a
diagonal norm and third-order closures stable in a block norm give global
third and fourth order. Test cases (their Section~7): the linear wave
equation, steady quasi-one-dimensional nozzle flow with the nonlinear Euler
equations, an unsteady two-dimensional Euler vortex, and a supersonic
hydrogen--air mixing layer with the Navier--Stokes equations; an eigenvalue
study shows the discrete operator stable for smooth and discontinuous data.
The appendices give the near-wall smoothness indicators, a pseudo-code for
the Euler equations, and the interpolation weights and dissipation
coefficients of both closures.

\section{What EXHALE uses today}

The WENO3 branch of \texttt{Reconstruction.f90} builds its smoothness
factors from Eq.~\eqref{eq:weights}: $\alpha_r = d_r(1 + \tau/(\epsilon +
\beta_r))$ with $\tau$ the squared second difference, $\beta_r$ the squared
one-sided differences, $d_0 = 2/3$, $d_1 = 1/3$, and the floor
$\epsilon = dr_j^2$, the local cell width squared, which carries the
$O(\Delta x^2)$ scaling of Eqs.~62 to 64 with the local width standing in for
the uniform $\Delta\xi$ and without the solution scale of Eq.~65. The code's
own note measures what that choice costs: nothing in a smooth region (the
face error and the convergence rate on the production grid are flat for
$\epsilon = (\kappa\, dr_j)^2$ over $\kappa$ from $10^{-1}$ to $10^{3}$,
$L_\infty$ $2.99\times10^{-8}$ at $N = 800$, rate 2.991, from the
\texttt{weno3\_reconstruction\_order} test), and something in the stencil
biasing, where on the WASP-121b reference output $\beta_r/\epsilon$ runs from
$8.4\times10^{3}$ at the base to $2.2\times10^{-6}$ at the top and a third of
the cells sit below $10^{-2}$, with no biasing left. So the code has the
accuracy result of the 2009a weights (Eq.~24), the design order at smooth
extrema included.

It does \emph{not} compute the right-hand side of Eq.~\eqref{eq:esweno}: no
$\mu_j$, no $\delta$, no dissipation flux (their Eq.~48). The same note says
so and gives three reasons the energy estimate would not carry over even if
it did:
\begin{enumerate}
\item The papers reconstruct the \emph{flux} by finite differences on a
  uniform grid (neither 2009 paper treats a non-uniform grid; the 2011 one
  makes the flux points non-uniform only near the boundaries and for a
  different reason). EXHALE reconstructs the primitive \emph{state} in a
  finite-volume sense on a radial grid whose spacing varies by a factor of
  about 20 across the domain, with volume-coordinate stencil coefficients of
  its own (\texttt{weno3\_geometry\_coefficients}), and hands the two states
  to an HLLC Riemann solver.
\item The estimate for a system holds in characteristic form. EXHALE
  reconstructs component by component, on primitive variables, which is the
  case the paper says it cannot bound.
\item The equations are the nonlinear Euler system with sources, for which
  the papers themselves claim no proof.
\end{enumerate}

\section{Assessment}

What ESWENO offers, in the order the papers establish it:
\begin{enumerate}
\item A proof, not an observation, that the semi-discrete operator does not
  create energy for linear hyperbolic problems, on any data, with no TVB
  assumption and no large-time-step constraint; and the converse statement
  that conventional WENO's symmetric part can have positive eigenvalues.
\item Design-order accuracy at smooth extrema with any number of vanishing
  derivatives, through the weights of Eq.~\eqref{eq:weights} and the
  $O(\Delta\xi^2)$ floor. This part EXHALE already has.
\item A systematic route to any order, at 15 to 20 percent more cost and the
  same stencil.
\item Boundary closures that keep the estimate on a finite domain, where
  conventional WENO closures lose either order or stability.
\end{enumerate}
What it does not offer: a proof for the nonlinear Euler equations, for a
component-wise reconstruction, or for a finite-volume reconstruction on a
stretched grid, which are the three properties of EXHALE's hydrodynamics.

For the questions open in this code the papers are therefore not a lever
yet. The base sound-wave oscillation and the marching limit cycle
(\texttt{docs/Update\_EXHALE\_stage1.pdf} items P44 and P54) were traced to
the boundary treatment and the time integration, not to the reconstruction
weights, and nothing measured on this code attributes an instability to a
positive eigenvalue of the reconstruction's symmetric part. Adopting the
dissipation term would first require restating the reconstruction as a
finite-difference flux operator in characteristic variables on the uniform
computational coordinate $\xi$ (the grid is a smooth map $r(\xi)$, so that
much is possible), and the finite-domain closures of \citet{Fisher2011}
would replace the ghost-cell boundary treatment at both ends. That is a
different hydrodynamic scheme, not an option to add to the present one, and
it would be judged as any such change is: by the certification of the
returned state and by the byte-identical regression on the cases whose
scheme is unchanged. No such work is planned.

\begin{thebibliography}{9}
\bibitem[Fisher et al.\ (2011)]{Fisher2011}
Fisher, T.~C., Carpenter, M.~H., Yamaleev, N.~K., \& Frankel, S.~H. 2011,
J. Comput. Phys., 230, 3727
\bibitem[Jiang \& Shu (1996)]{Jiang1996}
Jiang, G.-S., \& Shu, C.-W. 1996, J. Comput. Phys., 126, 202
\bibitem[Yamaleev \& Carpenter (2009a)]{Yamaleev2009a}
Yamaleev, N.~K., \& Carpenter, M.~H. 2009a, J. Comput. Phys., 228, 3025
\bibitem[Yamaleev \& Carpenter (2009b)]{Yamaleev2009b}
Yamaleev, N.~K., \& Carpenter, M.~H. 2009b, J. Comput. Phys., 228, 4248
\end{thebibliography}

\end{document}
